% Options for packages loaded elsewhere
\PassOptionsToPackage{unicode}{hyperref}
\PassOptionsToPackage{hyphens}{url}
%
\documentclass[
]{article}
\usepackage{amsmath,amssymb}
\usepackage{iftex}
\ifPDFTeX
  \usepackage[T1]{fontenc}
  \usepackage[utf8]{inputenc}
  \usepackage{textcomp} % provide euro and other symbols
\else % if luatex or xetex
  \usepackage{unicode-math} % this also loads fontspec
  \defaultfontfeatures{Scale=MatchLowercase}
  \defaultfontfeatures[\rmfamily]{Ligatures=TeX,Scale=1}
\fi
\usepackage{lmodern}
\ifPDFTeX\else
  % xetex/luatex font selection
  \setmainfont[]{FreeSerif}
\fi
% Use upquote if available, for straight quotes in verbatim environments
\IfFileExists{upquote.sty}{\usepackage{upquote}}{}
\IfFileExists{microtype.sty}{% use microtype if available
  \usepackage[]{microtype}
  \UseMicrotypeSet[protrusion]{basicmath} % disable protrusion for tt fonts
}{}
\makeatletter
\@ifundefined{KOMAClassName}{% if non-KOMA class
  \IfFileExists{parskip.sty}{%
    \usepackage{parskip}
  }{% else
    \setlength{\parindent}{0pt}
    \setlength{\parskip}{6pt plus 2pt minus 1pt}}
}{% if KOMA class
  \KOMAoptions{parskip=half}}
\makeatother
\usepackage{xcolor}
\usepackage[margin=1in]{geometry}
\setlength{\emergencystretch}{3em} % prevent overfull lines
\providecommand{\tightlist}{%
  \setlength{\itemsep}{0pt}\setlength{\parskip}{0pt}}
\setcounter{secnumdepth}{-\maxdimen} % remove section numbering
\ifLuaTeX
  \usepackage{selnolig}  % disable illegal ligatures
\fi
\IfFileExists{bookmark.sty}{\usepackage{bookmark}}{\usepackage{hyperref}}
\IfFileExists{xurl.sty}{\usepackage{xurl}}{} % add URL line breaks if available
\urlstyle{same}
\hypersetup{
  hidelinks,
  pdfcreator={LaTeX via pandoc}}

\title{Lemma F: far-field energies of the P2 error (step 4, T8)}
\author{}
\date{}

\begin{document}
\maketitle

\hypertarget{lemma-f-far-field-energies-of-the-p2-error-step-4-t8}{%
\section{Lemma F: far-field energies of the P2 error (step 4,
T8)}\label{lemma-f-far-field-energies-of-the-p2-error-step-4-t8}}

Status. §2 (Lemma F) and §4 (Lemma J) are proved here for the
\textbf{linear} error equation, with the modulation and residual forcing
included. The energies are half-curve energies; the far-end boundary
term is §3.1 (Lemma B and hypothesis (SE)). Lemma N3′ (§6.1) bounds the
nonlinear terms of the exact graph operator and the modulation terms 𝒬
at every level j ≤ 3, with x-dependent majorants on the certified far
boxes and without level 4 (\texttt{p2/t8se/n3sym.py},
\texttt{n3qsym.py}, \texttt{n3wcheck.py}, \texttt{n3far.py} →
\texttt{camp.py\ -\/-n3}). The algebraic core, the level-j identity
(I\_j), is machine-checked: \texttt{lfsym.py} (κ̃-coefficients, j ≤ 4)
and \texttt{cforms.py} (generic coefficients β, b₂, c₀, j ≤ 3) run an
Euler-operator test, and negative controls fail as they should. Every
constant is an explicit sup or integral of base quantities. Certified
(arb) values on the polished base, slabs 5--187: \texttt{cconst.py},
\texttt{merge.py}, \texttt{camp.py} → \texttt{constants/} (T8\_REPORT
§3). Float values: \texttt{fb.py}, \texttt{weights.py},
\texttt{lfnum.py}, \texttt{jcas.py}, \texttt{amp.py}. Inputs taken from
T5/T7 without being re-derived are tagged \textbf{{[}T7{]}}.

Notation: x is base arclength from the tip, y := x − X0 with X0 =
2.2490776515 (as in Lemma S), and c := ¼ − ã.

\hypertarget{setting}{%
\subsection{1. Setting}\label{setting}}

\textbf{1.1 The far-field equation.} T5 §2.1 (signs confirmed in T1
\texttt{signs.py}) gives e\_τ = 𝓛e − Σᵢxᵢgᵢ − (1−ϑ)r + 𝒩 + 𝒬, with 𝓛e =
−e'\,'\,'\,' + b₂e'\,' + b₁e' + b₀e and b₂ = −κ̃², b₁ = −5κ̃κ̃' − c q̃ − b̃t̃₁
+ ϖ, b₀ = c − 3κ̃`² − 4κ̃κ̃'`. Inserting ϖ = −κ̃κ̃' + Q + c(q̃ − x) + b̃t̃₁ (T1
§3) gives

b₁ = −c y + β, β := Q − cX0 − 6κ̃κ̃' + A, Q := ∫₀ˣκ̃`², b₀ = c + c₀, c₀ :=
−3κ̃'² − 4κ̃κ̃''.

Here A is the residual part of the transport (T6 §5: ϖ = ϖ* + A, A' =
κ̃r, A(0) = 0); the float study of T8\_REPORT §2 used A = 0. The terms
−cq̃ and −b̃t̃₁ cancel exactly in b₁, so b₁ involves neither the closure q̃
nor the turning t̃₁.

On the far field y ≥ y\_a ≥ 14.001 we have x ≥ 16.25, so ζ\_∞ = 1 and v
:= e' is T5's E\_I variable. Differentiating once:

\textbf{v\_τ = −v'\,'\,'\,' − c y v' + (βv + b₂v' + c₀e)' + f}, f :=
−x₁g₁' − x₂g₂' − x₃(gˢ₃ + gʳ)' − r' + (𝒩+𝒬)'. (1.1)

Here g₃ = Ṽ/n = gˢ₃ + gʳ + r/n with gˢ₃ := (cp̃ + b̃ñ)/n and gʳ := −κ̃'`/n,
by T5's definition r := Ṽ + κ̃'' − cp̃ − b̃ñ. So x₃g₃ + (1−ϑ)r = x₃(gˢ₃ +
gʳ) + r: the clock factor drops out of the residual forcing. The
derivation of (1.1) is ∂ₓ(𝓛e) = −v'\,'\,'\,' + (b₂v')' + (−cyv + βv)' +
cv + (c₀e)' = −v'\,'\,'\,' − cyv' + (βv + b₂v' + c₀e)'.

We write w\_m := ∂ₓᵐv for m ≥ 0 and w₋₁ := e. Lemma F treats f without
the nonlinear part (𝒩+𝒬)'; §6.1 adds 𝒩 and 𝒬 at levels 0--3.

\textbf{1.2 Parities.} The base is exactly symmetric about both ends
(T1: κ̃' = κ̃'\,'\,' = 0 at x = 0 and x = ℓ̃). We view the half curve as
one half of the closed doubled curve x ∈ {[}0, 2ℓ̃{]}. About the far end
x = ℓ̃:

\begin{itemize}
\tightlist
\item
  κ̃, b₂, c₀ and e are even, and w\_m has parity (−1)\^{}\{m+1\};
\item
  β\^{}\{(m)\} has parity (−1)\^{}\{m+1\} for m ≥ 1;
\item
  b₁ + ℓ̃\_τ is odd up to the constant D := b₁(ℓ̃) + ℓ̃\_τ. This holds
  because κ̃κ̃', Q − Q(ℓ̃) and y − y\_ℓ are odd, with y\_ℓ := ℓ̃ − X0. With
  T1's length equation ℓ̃\_τ = cℓ̃ − Q(ℓ̃) + ε\_ℓ and b₁(ℓ̃) = Q(ℓ̃) − cℓ̃ +
  A(ℓ̃) one gets \textbf{D = A(ℓ̃) + ε\_ℓ}, \textbar D\textbar{} ≤ 6.9e-13
  (T6: sup\textbar A(ℓ̃)\textbar{} ≤ 6.7e-13, sup\textbar ε\_ℓ\textbar{}
  ≤ 2.3e-14). T6's scalar D\_end = ϖ(ℓ̃) + ℓ̃\_τ equals D + cq̃(ℓ̃) +
  b̃t̃₁(ℓ̃); it is not the transport coefficient (the closure terms cancel
  in b₁, §1.1).
\end{itemize}

\textbf{1.3 Weights.} For each level j ∈ \{0..4\} take a ramp {[}y\_a,j,
y\_b,j{]} and set T\_j(x,τ) := 28 η\_j(y)² Y\_ℓ(x,τ)\^{}\{2j−2\}
Ω(Y\_ℓ/Y\_h(τ)), where: - η\_j = S(t), S(u) = u⁴(35 − 84u + 70u² −
20u³), t = ln(y/y\_a)/ln(y\_b/y\_a), clipped to {[}0,1{]}. The T5 design
is y\_a, y\_b = 14, 28 for all j (the Lemma S ramp). - Y\_ℓ := y\_ℓ
Φ(y/y\_ℓ) is a far-end flattening. Φ(s) = s for s ≤ ½, and Φ' = 1 −
ψ(2s−1) with ψ(u) = u⁵(21/4 − 49u/8 + 15u²/8). Then Φ' is odd about s =
1, so Y\_ℓ is even about ℓ̃ on the doubled curve, and Y\_ℓ = y on the
half of the curve nearer the tip. - Ω ≡ 1 for the \emph{global}
energies. The \emph{localized} energies of Remark 6.3 use an outer
factor that follows the characteristics.

The properties used are: (W1) T\_j ∈ W\^{}\{4,∞\}, since η² is C³ at
y\_b with T'\,'\,'\,' bounded; (W2) T\_j = 0 for y ≤ y\_a and T\_j ≍ (y
− y\_a)⁸ there, so T'²/T and T'\,'²/T are bounded; (W3) T\_j is even
about ℓ̃. (The jets are checked against finite differences in
\texttt{weights.py}; the only mismatch is at the C³ point y = y\_b.)

Regions: the ramp R\_j := \{y\_a,j \textless{} y \textless{} y\_b,j\}
and the core C\_j := \{y ≥ y\_b,j\}, the latter including the flattened
far end.

\textbf{1.4 Energies and data.} The energies are the half-curve
integrals F\_j := ∫\_\{x*\}\^{}\{ℓ̃(τ)\} T\_j w\_j², the dissipations
D\_j := ∫T\_j w\_\{j+2\}², and N\_\{j+1\}² := ∫T\_j w\_\{j+1\}² (same
domain). The ramp data are the E\_I norms of T5 restricted to the ramp,
V\_m := (∫\_\{R\} W (∂ᵐv)²)\^{}\{1/2\} ≤ T5's V\_m, with W :=
(1+x²)\^{}\{−2\}. The cost of a stronger weight is ∫T\_j w\_j² vs ∫W
w\_j², i.e.~y\^{}\{2j+2\} on the ramp. This conversion factor drives the
amplification (T8\_REPORT §2).

\textbf{1.5 Regularity.} κ̃ ∈ C⁴ with κ̃\^{}\{(5)\} ∈ L\^{}∞ (T1), so e ∈
W\^{}\{7,∞\} and v ∈ W\^{}\{6,∞\} on the far field (T5 §2.1). The
identity (I\_j) below involves w\_\{≤ j+2\} and κ̃\^{}\{(≤5)\} only. It
is proved for smooth data and extends by density, because every term is
continuous in the norms that are available. The time derivative is taken
a.e.; F\_j is absolutely continuous.

\hypertarget{statement}{%
\subsection{2. Statement}\label{statement}}

\textbf{Lemma F.} Let j ∈ \{0,\ldots,4\}, let the weights satisfy
(W1)--(W3), and let the ramps satisfy y\_a,i \textless{} y\_b,j for i
\textless{} j, so that T\_i \textgreater{} 0 on C\_j. Let F\_j be the
half-curve energies of §1.4 and assume the symmetric-end hypothesis (SE)
of §3.1. Along the linear part of (1.1), for a.e. τ:

F\_j' ≤ −λ\_j F\_j − θ\_j D\_j + 2√F\_j g\_j + h\_j, (F)

g\_j = Σ\_\{i\textless j\} ½c\_ji √F\_i + \textbar x\textbar γ\_j +
½a\_j Ȳ, h\_j = Σ\_\{i\textless j\} d\_ji F\_i + 𝓘\_j + e\_j Ȳ² +
δ\_s⁻¹(𝓡\_j + \textbar x\textbar γʳ\_j)² + 𝓡̂\_j²,

with the ramp data

𝓘\_j = k\_jj V\_j² + Σ\_\{i\textless j\}(k\_ji V\_i² + r\_ji V\_j V\_i)
+ re\_j V\_j Ȳ (standard ramp), 𝓘\_j = kD\_j D\_\{j−2\} +
Σ\_\{i\textless j\} sD\_ji (D\_\{j−2\}F\_i)\^{}\{1/2\} + \ldots{}
(nested ramp: R\_j ⊂ C\_\{j−2\}).

The constants are defined as follows. Fix a = 0.01, ν = 10⁻³, δ\_P ∈ (0,
1.5{]}, δ\_s ∈ (0, ½{]} and δ₆ = 0.1. Let n\_j be the number of
(j+1,·)-pairs in (I\_j), plus 1 for the forcing, recall P\_j, q\_j,
p\_ji from §3.2, and let ε\_j := \textbar D\textbar/m\_j be the far-end
constant of Lemma B (§3.1).

\begin{itemize}
\tightlist
\item
  q̂\_j := ∂\_τT\_j + q\_j + ½P\_j'\,' + P\_j²/(4δ\_P T\_j) + (a n\_j/4 +
  ε\_j/2)(T\_j'`)⁺ + ν\textbar T\_j''\textbar.
\item
  ΔF\_j := a n\_j/8 + ε\_j(2/y\_ℓ + 5/4), σ\_j := c + sup\_\{C\_j\}
  q̂\_j/T\_j + ΔF\_j and λ\_j := c − σ\_j.
\item
  θ\_j := 2 − δ\_P − a n\_j/2 − δ₆{[}j=4{]} − δ\_s − ε\_j. (The design
  opt4n has δ\_P = 1.5 at j = 3; it needs δ\_s \textless{} ½ there, and
  the certified run uses δ\_s = ¼.)
\item
  k\_jj := sup\_\{R\_j\} (q̂\_j + (λ\_j + ΔF\_j) T\_j)⁺/W. In the nested
  case, kD\_j := sup\_\{R\_j\} (q̂\_j + (λ\_j + ΔF\_j)T\_j)⁺/T\_\{j−2\}.
  (Earlier versions omitted ΔF\_jT\_j here; it is the ramp share of the
  F\_j-part of the Young steps, §3.4.)
\item
  c\_ji := sup\_\{C\_j\} \textbar p\_ji\textbar/√(T\_jT\_i). The
  e-coupling adds 2A₂ to c\_j0; A₂ would suffice.
\item
  r\_ji := sup\_\{R\_j\}\textbar p\_ji\textbar/W. In the nested case,
  sD\_ji := sup\_\{R\_j\}\textbar p\_ji\textbar/√(T\_\{j−2\}T\_i).
\item
  (j+1,0)-pairs p: d\_j0 += sup\_\{C\_j\} p²/(T\_jT\_0)/(2a) and k\_j0
  += sup\_\{R\_j\} p²/(T\_jW)/(2a).
\item
  (j,e)-pairs p: a\_j := ε̂₀(∫\_\{C\_j\}p²/T\_j)\^{}\{1/2\}, A₂ :=
  (∫\_\{C\_j\}(p²/T\_j)H²)\^{}\{1/2\}, and re\_j :=
  (sup\_\{R\_j\}\textbar p\textbar/W ·
  ∫\_\{R\_j\}\textbar p\textbar ê²)\^{}\{1/2\}.
\item
  (j+1,e)- and (j+2,e)-pairs p: e\_j += {[}2ε̂₀²∫\_\{C\_j\}p²/T\_j +
  ∫\_\{R\_j\}(p²/T\_j)ê²{]}/κ and d\_j0 += 2∫\_\{C\_j\}(p²/T\_j)H²/κ,
  with κ = 2a for (j+1,e) and κ = 4δ₆ for (j+2,e).
\item
  γ\_j := (Σᵢ∫T\_j(∂\^{}\{j+1\}gᵢˢ)²)\^{}\{1/2\} and γʳ\_j :=
  (∫T\_j(∂\^{}\{j−1\}gʳ)²)\^{}\{1/2\}, where gˢ := (p̃/‖χp̃‖\_ρ,
  ñ/‖χñ‖\_ρ, (cp̃+b̃ñ)/n) and gʳ := −κ̃'`/n with n = ‖χṼ‖\_ρ (T5
  normalisation). For j = 0, ∂\^{}\{−1\}gʳ := −κ̃'/n.
\item
  𝓡\_j := (∫T\_j(∂\^{}\{j−1\}r)²)\^{}\{1/2\}, which is T5's R\^{}F\_j up
  to the Y\_ℓ flattening, with ∂\^{}\{−1\}r := −∫\_x\^{}\{ℓ̃\}r.
\item
  𝓡̂\_j² := (8/a)∫(T\_j'²/T\_j)ρ\_j² +
  ν⁻¹∫\textbar T\_j'\,'\textbar ρ\_j², with ρ\_j := ∂\^{}\{j−1\}(r +
  x₃gʳ).
\end{itemize}

The e-bound uses three functions:

\begin{itemize}
\tightlist
\item
  ê(x) := K\_e + Ĩ(x)/√μ\_e, with Ĩ(x)² := ∫\_\{x*\}\^{}x W\^{}\{−1\}
  and \textbar e(x*)\textbar{} ≤ K\_eȲ \textbf{{[}T7{]}};
\item
  ε̂₀ := ê(x₁), where x₁ is the start of C₀;
\item
  H(x)² := ∫\_\{x₁\}\^{}x T₀\^{}\{−1\} on C₀, with H := 0 below x₁.
\end{itemize}

\textbf{Validity.} All sups are finite by (W2). If λ\_j \textgreater{} 0
and θ\_j \textgreater{} 0, then (F) holds with these constants (any λ ≤
λ\_j may be used in (F) and in k\_jj). In the nested case, R\_j ⊂
C\_\{j−2\} and T\_i \textgreater{} 0 on R\_j are required.

\textbf{Hand-off coverage.} The window energies of Lemma S (cutoff
S(ln(y/14)/ln 2) on {[}14,28{]}) are dominated by F\_j whenever η\_j ≥
S(ln(y/14)/ln 2), for instance for every log-ramp {[}14, y\_b{]} with
y\_b ≤ 28.

\hypertarget{proof}{%
\subsection{3. Proof}\label{proof}}

\textbf{3.1 Half-curve energies and the far-end boundary term.} T\_j
vanishes to order 8 at y\_a (W2), so nothing happens at x*. At the
moving far end,

F\_j' = ∫\_\{x*\}\^{}\{ℓ̃\}(∂\_τT\_j)w\_j² + 2∫\_\{x*\}\^{}\{ℓ̃\}T\_j w\_j
∂\_τw\_j + ℓ̃\_τT\_j(ℓ̃)w\_j(ℓ̃)².

Insert (I\_j); its total-derivative part 𝔉\_j leaves 𝔉\_j(ℓ̃). The flux
of the transport term 2T\_jw\_j·b₁w\_\{j+1\} = (T\_jb₁w\_j²)' −
(T\_jb₁)'w\_j² is T\_jb₁w\_j². Hence

F\_j' = ∫\_\{x*\}\^{}\{ℓ̃\}{[}(∂\_τT\_j)w\_j² + density of (I\_j){]} +
𝔅\^{}D\_j + 𝔅\^{}S\_j, 𝔅\^{}D\_j := D·T\_j(ℓ̃)w\_j(ℓ̃)², D = b₁(ℓ̃) + ℓ̃\_τ
= A(ℓ̃) + ε\_ℓ (§1.2),

and 𝔅\^{}S\_j is the sum of the fluxes at ℓ̃ of the fourth-order, b₂-,
c₀-, forcing and X\_j-terms. 𝔅\^{}S\_j contains no b₁; it is a bilinear
form in the boundary jets (e, w₀, \ldots, w\_\{j+3\})(ℓ̃) with
coefficients from T\_j\^{}\{(≤3)\}(ℓ̃) (T\_j' = T\_j'\,'\,' = 0 there by
(W3)) and the coefficient jets at ℓ̃.

\textbf{Lemma B (transport defect).} With m\_j :=
(Φ(½)/Φ(1))\^{}\{2j−2\} for j ≥ 2, m₀ = m₁ = 1 (Φ(½) = ½, Φ(1) =
0.8828125) and ε\_j := \textbar D\textbar/m\_j,

\textbar 𝔅\^{}D\_j\textbar{} ≤ ε\_j{[}(2/y\_ℓ + 5/4)F\_j + D\_j +
½∫(T\_j'\,')⁺w\_j²{]}.

\emph{Proof.} P := \{y\_ℓ/2 ≤ y ≤ y\_ℓ\} has length d = y\_ℓ/2 and lies
in C\_j. On P, T\_j ≥ m\_jT\_j(ℓ̃): T₀ = 28Y\_ℓ⁻² decreases, T₁ = 28 is
constant, T\_j increases for j ≥ 2. For f = w\_j, the trace inequality
f(ℓ̃)² ≤ d⁻¹∫\_P f² + 2‖f‖\_\{L²(P)\}‖f'‖\_\{L²(P)\} (average f(ℓ̃)² =
f(z)² + ∫\_z\^{}ℓ̃(f²)' over z ∈ P) gives T\_j(ℓ̃)w\_j(ℓ̃)² ≤ m\_j⁻¹(F\_j/d
+ 2√F\_jN\_\{j+1\}) ≤ m\_j⁻¹((2/y\_ℓ + 1)F\_j + N²\_\{j+1\}). Finally
N²\_\{j+1\} ≤ D\_j + ¼F\_j + ½∫(T\_j'\,')⁺w\_j² (S3). ∎

Lemma B is what the constants ΔF\_j, θ\_j and q̂\_j of §2 contain. With
\textbar D\textbar{} ≤ 6.9e-13, ε\_j ≤ 2.1e-11: negligible, but included
in the certified run.

\textbf{(SE) Symmetric end.} 𝔅\^{}S\_j ≤ 0.

(SE) holds with 𝔅\^{}S\_j = 0 when the error, the base coefficients and
the residual are symmetric about the far-end normal line. This is the
closed-curve argument: write 2F\_j as an integral over the doubled far
field {[}x*, 2ℓ̃ − x*{]}, on which (1.1) holds with the even/odd
continuations of the coefficients (§1.2). Under x ↦ 2ℓ̃ − x every density
is even by §1.2 and (W3), and the transport/∂\_τT pair is
mirror-symmetric (b₁(2ℓ̃−x) = −b₁(x) − 2ℓ̃\_τ, ∂\_τ{[}T(2ℓ̃−x){]} =
(∂\_τT)(2ℓ̃−x) + 2ℓ̃\_τT'(2ℓ̃−x)), so no flux appears at ℓ̃. The evenness of
T\_j is what makes the fourth-order fluxes vanish; with the unflattened
power weight, 2T'w\_jw\_j'\,' + T'`'w\_j² would survive at ℓ̃ for odd j,
and a positive term 2\textbar T'\textbar w\_j'² for j = 0.

\emph{Remark (far end off the axis).} The base is symmetric about its
own far-end normal line N(τ), since κ̃ is exactly even about ℓ̃. The true
curve is symmetric about the axis. By T6, N is off the axis by
\textbar q̃(ℓ̃)\textbar{} ≤ 4.6e-9 and tilted by \textbar t̃₁(ℓ̃)\textbar{}
≤ 2.4e-13. The offset grows like ℓ̃ (q̃(ℓ̃)/ℓ̃ ≈ 1.3e-14 on slabs ≥ 100),
and the turning defect θ̃(ℓ̃) + 3π/2 ≤ 1.65e-14 is constant. This is the
signature of a fixed physical rigid rotation of about 1.6e-14 rad. So
(SE) fails only through this rigid mode; by §1.1 the offset does
\emph{not} enter the transport term. There are two ways to restore (SE):

\begin{itemize}
\item
  \begin{enumerate}
  \def\labelenumi{(\alph{enumi})}
  \tightlist
  \item
    Compare the true curve with the base rotated by that angle.
    Rotations about the origin are symmetries of the rescaled flow up to
    the axial drift, which changes r by at most
    \textbar b̃\textbar·1.7e-14 ≤ 1.1e-15 (T6 would enclose this). This
    is a choice in P2's representation, for T5/T9; T4 must check the
    start tolerance, since the tip moves by about 4e-14.
  \end{enumerate}
\item
  \begin{enumerate}
  \def\labelenumi{(\alph{enumi})}
  \setcounter{enumi}{1}
  \tightlist
  \item
    Bound 𝔅\^{}S\_j directly, as (jets of the rigid mode at ℓ̃, of size
    1.7e-14·ℓ̃(19.4/ℓ̃)\^{}m) × (traces of w\_\{≤j+3\} at ℓ̃). This needs
    pointwise far-end control of w\_\{j+3\}, which the level-j energies
    do not give.
  \end{enumerate}
\end{itemize}

The mode's own energies are ≤ 1e-20 relative to F\_j, so the effect is
nil; but (SE) is an input of Lemma F, not a consequence.

\textbf{3.2 The identity (I\_j)} (machine-checked: \texttt{lfsym.py},
output \texttt{lf\textbackslash{}\_forms.pkl}). With T = T\_j,

2T w\_j ∂ʲ{[}v\_τ − f{]} ≐ −2T w\_\{j+2\}² + P\_j w\_\{j+1\}² + q\_j
w\_j² + Σ\_\{i=−1\}\^{}\{j−1\} p\_ji w\_j w\_i + X\_j,

where ≐ means equality modulo a total x-derivative, and

\begin{itemize}
\tightlist
\item
  P\_j = 4T'\,' + 2κ̃²T. The second summand is the anti-diffusive −2b₂T
  term.
\item
  q\_j = −T'\,'\,'\,' + (cy − β)T' + (1−2j)cT + (2j+1)β'T + (Tb₂)'\,' −
  (j+1)(Tb₂')' + j(j+1)Tb₂'\,' + 2c₀T.
\item
  p\_ji = 2C(j+1,j+1−i)Tβ\^{}\{(j+1−i)\}·{[}i ≥ 0, j+1−i ≥ 2{]} +
  2C(j+1,j+2−i)Tb₂\^{}\{(j+2−i)\}·{[}1 ≤ i ≤ j−1{]} +
  2C(j+1,j−i)Tc₀\^{}\{(j−i)\}·{[}i ≤ j−1{]}.
\item
  X\_j keeps every coefficient derivative at κ̃\^{}\{(≤5)\}. It replaces
  three high-order terms by integrations by parts onto (Tw\_jw\_i):

  \begin{itemize}
  \tightlist
  \item
    j=3: the c₀\^{}\{(4)\}-term 2Tc₀\^{}\{(4)\}w₃e becomes
    −2c₀'\,'\,'(T'w₃e + Tw₄e + Tw₃w₀).
  \item
    j=4, β\^{}\{(5)\} and c₀\^{}\{(4)\} terms: these become
    −(2β\^{}\{(4)\} + 10c₀'\,'\,')(T'w₄w₀ + Tw₅w₀ + Tw₄w₁).
  \item
    j=4, c₀\^{}\{(5)\}-term: this becomes 2c₀'\,'`(T''w₄e + 2T'w₅e +
    2T'w₄w₀ + Tw₆e + 2Tw₅w₀ + Tw₄w₁).
  \end{itemize}
\end{itemize}

Derivation:

\begin{itemize}
\tightlist
\item
  From −2Tw\_jw\_j'\,'\,'`, three integrations by parts give
  −2Tw\_\{j+2\}² + 4T'`w\_\{j+1\}² − T''\,'\,'w\_j².
\item
  From ∂ʲ(−cyw₁) = −cyw\_\{j+1\} − jcw\_j: the term c(yT' +
  (1−2j)T)w\_j².
\item
  The β-, b₂- and c₀-terms are expanded by Leibniz. The β and b₂ terms
  with m = 0, 1 (and m = 2 for b₂) are folded into q\_j and P\_j by one
  or two integrations by parts.
\end{itemize}

The Euler test checks the whole identity for arbitrary T, κ̃ and Q with
Q' = κ̃'². It reports κ̃-derivatives ≤ 5 in all coefficients. Pairs
present: j=0: (0,e); j=1: (1,e),(1,0); j=2: (2,e),(2,0),(2,1); j=3:
(3,e),(3,0..2),(4,e); j=4: (4,e),(4,0..3),(5,e),(5,0),(6,e).

\textbf{3.3 Bounds.}

\begin{itemize}
\tightlist
\item
  \textbf{(S1)} Exactly, ∫P w\_\{j+1\}² = −∫P w\_j w\_\{j+2\} +
  ½∫P'\,'w\_j². Then \textbar∫Pw\_jw\_\{j+2\}\textbar{} ≤ δ\_P D\_j +
  (4δ\_P)⁻¹∫(P²/T)w\_j².
\item
  \textbf{(S2) (j,i), 0 ≤ i \textless{} j.} On C\_j the weights are
  regular: ∫\textbar p\textbar\textbar w\_jw\_i\textbar{} ≤
  c\_ji√(F\_jF\_i). On R\_j use the ramp data: ≤ r\_ji V\_jV\_i. In the
  nested case use instead ≤
  (∫\_\{R\_j\}T\_\{j−2\}w\_j²)\^{}\{1/2\}(∫(p²/T\_\{j−2\})w\_i²)\^{}\{1/2\}
  ≤ sD\_ji√(D\_\{j−2\}F\_i).
\item
  \textbf{(S3) (j+1,i), i ∈ \{0,e\}.} First N\_\{j+1\}² =
  −∫Tw\_jw\_\{j+2\} + ½∫T'`w\_j² ≤ D\_j + ¼F\_j + ½∫(T'')⁺w\_j². Then
  \textbar∫pw\_\{j+1\}w\_i\textbar{} ≤ N\_\{j+1\}M\_i ≤ (a/2)N\_\{j+1\}²
  + M\_i²/(2a), where M\_i² = ∫(p²/T)w\_i² is split over C\_j (giving
  F₀) and R\_j (giving V₀ or ê).
\item
  \textbf{(S4) (j+2,e).} \textbar∫pw\_\{j+2\}e\textbar{} ≤ δ₆D\_j +
  (4δ₆)⁻¹∫(p²/T)e².
\item
  \textbf{(S5) e.} Write e(x) = e(x*) + ∫\_\{x*\}\^{}x v. Then
  \textbar e\textbar{} ≤ ê(x)Ȳ on R₀, using μ\_eV₀² ≤ Λ\_e ≤ Ȳ². On C₀,
  \textbar e\textbar{} ≤ ε̂₀Ȳ + H(x)√F₀, by Cauchy--Schwarz with T₀ and
  because ê is increasing. Hence
  ∫\_\{C\_j\}\textbar p\textbar\textbar w\_j\textbar\textbar e\textbar{}
  ≤ √F\_j(a\_jȲ + A₂√F₀) by Minkowski, and
  ∫\_\{R\_j\}\textbar p\textbar\textbar w\_j\textbar\textbar e\textbar{}
  ≤ re\_jV\_jȲ. For j ≥ 1 the A₂-part goes into c\_j0. For j = 0 it is a
  self-coupling, and λ₀ is lowered by 2A₂ (≤ 10⁻⁸; the certified run
  does this).
\item
  \textbf{(S6) Forcing.}

  \begin{itemize}
  \tightlist
  \item
    Residual and rough modulation: two integrations by parts give
    −2∫(Tw\_j)`'ρ\_j = −2∫(T'`w\_j + 2T'w\_\{j+1\} + Tw\_\{j+2\})ρ\_j
    (their fluxes at ℓ̃ belong to 𝔅\^{}S\_j). With ‖ρ\_j‖\_T ≤ 𝓡\_j +
    \textbar x\textbar γʳ\_j, this is ≤ δ\_sD\_j + δ\_s⁻¹(𝓡\_j +
    \textbar x\textbar γʳ\_j)² + (a/2)N\_\{j+1\}² + 𝓡̂\_j² +
    ν∫\textbar T''\textbar w\_j². (For j=0, ∂\^{}\{−1\}r = S₁ − S₁(ℓ̃)
    vanishes at ℓ̃.)
  \item
    Smooth modulation: Cauchy--Schwarz in x and in i gives ≤
    2\textbar x\textbar√F\_jγ\_j. For this smooth forcing the
    two-integrations route is far worse: it pays the y⁴ weight mismatch
    between D\_j and F\_j.
  \end{itemize}
\end{itemize}

\textbf{3.4 Collecting.}

\begin{itemize}
\tightlist
\item
  All w\_j²-terms are ∫q̂\_jw\_j² + ΔF\_jF\_j, where ΔF\_j is the F\_j
  part of the N\_\{j+1\}² bounds in (S3), (S6) and Lemma B. On C\_j they
  are ≤ (sup\_\{C\_j\}q̂\_j/T\_j + ΔF\_j)∫\_\{C\_j\}T\_jw\_j² =
  −λ\_j∫\_\{C\_j\}T\_jw\_j². On R\_j, write ∫\_\{R\_j\}(q̂\_j +
  ΔF\_jT\_j)w\_j² = −λ\_j∫\_\{R\_j\}T\_jw\_j² + ∫\_\{R\_j\}(q̂\_j + (λ\_j
  + ΔF\_j)T\_j)w\_j²; the last integral is ≤ k\_jjV\_j². In the nested
  case it is ≤ kD\_jD\_\{j−2\}, using T\_\{j−2\} ≡
  T\_\{j−2\}\textbar\_\{core\} on R\_j.
\item
  The D\_j-coefficient is −2 + δ\_P + an\_j/2 + δ₆{[}j=4{]} + δ\_s +
  ε\_j = −θ\_j.
\end{itemize}

Together with 𝔅\^{}D\_j (Lemma B) and 𝔅\^{}S\_j ≤ 0 (SE), this proves
(F). ∎

\hypertarget{lemma-j-time-weighted-ramp-data}{%
\subsection{4. Lemma J (time-weighted ramp
data)}\label{lemma-j-time-weighted-ramp-data}}

\textbf{Hypotheses.} T5's closure is run with a reserved fraction κ\_𝔇
of its dissipation surplus. Concretely:

\begin{itemize}
\tightlist
\item
  ½Λ\_e' ≤ −ϱΛ\_e − κ\_𝔇𝔇 + a\_cY + Ξ'², where 𝔇 = (θ₀/4)B² + (ε/8)D₄² +
  (μ\_e/8)V₂² + (μ\_eε\_I/8)V₄² and Ξ'² = Ξ²/(1−κ\_𝔇).
\item
  Ȳ' = −ϱȲ + Φ with Φ = a\_c + √ϱΞ', √Λ\_e ≤ Ȳ ≤ 1 and Ȳ ≥ Ξ'/√ϱ
  \textbf{{[}T7{]}}.
\end{itemize}

As written in T5 (Y), the surplus is consumed by the residual Young
steps, so κ\_𝔇 = 0 there; we use κ\_𝔇 = ½.

We also use:

\begin{itemize}
\tightlist
\item
  (I0) and (I2) of T5 §2.4 with \textbar x\textbar{} ≤ C\_xA;
\item
  the interpolations W₁² ≤ AB + c₁A², W₃² ≤ BD₄ + c₁B², V₁² ≤ V₀V₂ + c₁W
  V₀², V₃² ≤ V₂V₄ + c₁W V₂² \textbf{{[}T7/S3{]}}.
\end{itemize}

Let 0 \textless{} λ \textless{} 2 min(μ₁, μ₂), ω(s) := e\^{}\{−λ(τ−s)\},
and J\_X := ∫\_\{τ0\}\^{}τ ωX².

\textbf{Conclusion.}

\begin{enumerate}
\def\labelenumi{\arabic{enumi}.}
\tightlist
\item
  ∫ω𝔇 ≤ 𝔅 := κ\_𝔇⁻¹{[}½ + (ϱ̄ − λ/2)⁺/λ{]}.
\item
  J\_A ≤ 1/λ, J\_B ≤ 4𝔅/θ₀, J\_\{D4\} ≤ 8𝔅/ε, J\_\{W1\} ≤ √(J\_AJ\_B) +
  c₁J\_A, J\_\{W3\} ≤ √(J\_BJ\_\{D4\}) + c₁J\_B.
\item
  (μ₁ − λ/2)J\_\{V0\} + ½J\_\{V2\} ≤ K + C\_xγ₀√(J\_AJ\_\{V0\}), with K
  := ½ω(τ0)/μ\_e + a\_fJ\_\{W1\} + b\_fJ\_B.
\item
  J\_\{V4\} ≤ 2ω(τ0)/(μ\_eε\_I) + M(λ)(c₃√J\_\{D4\} + c₂√J\_\{W3\} +
  c₁ᴵ√J\_B + c₀ᴵ√J\_\{W1\} + C\_xγ₂√J\_A)², with M(λ) := min\_κ
  4{[}(1+Ψ₁/κ)² + (Ψ₁κ + 2Ψ₁√c₁W + 2c₁W)²/(4(μ₂ − λ/2)){]}.
\item
  J\_\{V1\} ≤ √(J\_\{V0\}J\_\{V2\}) + c₁W J\_\{V0\} and J\_\{V3\} ≤
  √(J\_\{V2\}J\_\{V4\}) + c₁W J\_\{V2\}. The Λ-routes J\_\{V2\} ≤
  8𝔅/μ\_e and J\_\{V4\} ≤ 8𝔅/(μ\_eε\_I) are also valid.
\item
  Residual terms. Suppose the majorant satisfies Ȳ ≥ Ȳ\_* (the certified
  run uses Ȳ\_* = 10⁻¹²; the budget is about 2.3·10⁻¹²). The term 𝓡₁(V₂
  + 2Ψ₁V₁ + Ψ₂V₀) of (I0) is bounded by δ(V₀² + V₂²) + 𝓡₁²C\_res/(4δ),
  where C\_res = (1+Ψ₁)² + (Ψ₁(1 + 2√c₁W) + Ψ₂)² and Ψ₂ := 2c₁W (this
  uses V₁ ≤ ½(V₀ + V₂) + √c₁W V₀). Hence item 3 holds with μ₁ → μ₁ − δ,
  ½J\_\{V2\} → (½ − δ)J\_\{V2\}, and K += C\_res(𝓡̄₁/Ȳ\_*)²/(4δλ). In
  (I2), 𝓡ᴵ\_s enters X additively; it adds 𝓡̄ᴵ\_s/(Ȳ\_*√λ) to √∫ωX². The
  run uses δ = 10⁻³.
\end{enumerate}

\textbf{Proof.}

\begin{enumerate}
\def\labelenumi{\arabic{enumi}.}
\tightlist
\item
  Multiply the closure by ω and integrate, using ½∫ωΛ' = ½{[}ωΛ{]} −
  (λ/2)∫ωΛ. Since Y ≤ Ȳ and Ξ'² ≤ √ϱΞ'Ȳ, we have a\_cY + Ξ'² ≤ ΦȲ. Write
  ∫ωΦȲ = ∫ω(½(Ȳ²)' + ϱȲ²) = ½Ȳ(τ)² − ½ω(τ0)Ȳ(τ0)² + ∫ω(ϱ − λ/2)Ȳ². This
  gives κ\_𝔇∫ω𝔇 ≤ ½Ȳ(τ)² + ∫ω(ϱ − λ/2)(Ȳ² − Λ\_e). The last integral is
  ≤ (ϱ̄ − λ/2)⁺/λ.
\item
  The pointwise bound Λ\_e ≤ Ȳ² ≤ 1 gives J\_A, and item 1 gives J\_B
  and J\_\{D4\}. The remaining bounds follow from the interpolations and
  Cauchy--Schwarz in time.
\item
  Multiply (I0) by ω.
\item
  Multiply (I2) by ω. Use V₃ ≤ √(V₂V₄) + √c₁W V₂ ≤ V₄/(2κ) + (κ/2 +
  √c₁W)V₂, then Young: XαV₄ ≤ ¼V₄² + α²X² and Xβ̂V₂ ≤ (μ₂ − λ/2)V₂² +
  β̂²X²/(4(μ₂ − λ/2)). Finally apply Minkowski in time to √∫ωX². ∎
\end{enumerate}

For the T7 inputs \textbf{{[}T7{]}} (S1 reference values / T5 defaults):
μ\_e = 0.02, ε = 10⁻⁵, ε\_I = 3·10⁻⁵, θ₀ = 0.02, ϱ̄ = 0.03, c₁ = 0.5804,
a\_f = 0.0696, b\_f = 0.0394, μ₁ = 0.1552, μ₂ = 0.6555, c\^{}I\_\{0..3\}
= (0.425, 0.424, 0.526, 0.442), c₁W = 0.8477, Ψ₁ = 0.39, C\_x = 4, γ₀ =
0.05, γ₂ = 0.01, κ\_𝔇 = ½. With λ = 0.244 this gives, per unit Ȳ²:

\begin{itemize}
\tightlist
\item
  J₀..J₄ = (205, 255, 31.7, 1.24·10⁴, 4.82·10⁶);
\item
  73\% of J₄ is the c₃²J\_\{D4\} = c₃²·8𝔅/ε term, i.e.~the 1/√ε loss.
\end{itemize}

\hypertarget{corollary-c-comparison-and-hand-off}{%
\subsection{5. Corollary C (comparison and
hand-off)}\label{corollary-c-comparison-and-hand-off}}

\textbf{C1 (supersolution).} Suppose F' ≤ −λF + 2g√F + h with λ, g, h ≥
0. Define H := e\^{}\{−∫λ\}F(τ0) + ∫e\^{}\{−∫\_s\^{}τλ\}h and G :=
∫e\^{}\{−½∫\_s\^{}τλ\}g. Then √F ≤ √H + G.

Proof. Z := (√H + G)² satisfies Z' = −λZ + 2g√Z + h + Gh/√H ≥ −λZ + 2g√Z
+ h. On any interval where F \textgreater{} Z \textgreater{} 0 we get (F
− Z)' ≤ (−λ + g/√Z)(F − Z). Grönwall from the last time with F ≤ Z gives
F ≤ Z; add δ \textgreater{} 0 to F(τ0) and let δ → 0. ∎

\textbf{C2 (triangular system).} Apply C1 level by level, with √F\_i ≤
√H\_i + G\_i for i \textless{} j inside g\_j and h\_j, and with ∫ωh\_j
bounded through Lemma J:

\begin{itemize}
\tightlist
\item
  ∫ω𝓘\_j ≤ Σ k\_jm J\_m + Σ r\_ji √(J\_jJ\_i) + re\_j √(J\_j/λ);
\item
  in the nested case, ∫ωD\_\{j−2\} ≤
  θ\_\{j−2\}⁻¹{[}H\_\{j−2\}\^{}\{(λ)\} + 2
  sup(g\_\{j−2\}√F\_\{j−2\})/λ{]}. Here ω and H\_\{j−2\}\^{}\{(λ)\} use
  the rate λ = min(λ\_j, λ\_\{j−2\}), since the term (λ −
  λ\_\{j−2\})⁺∫ωF\_\{j−2\} must vanish.
\end{itemize}

\texttt{amp.py} implements this with per-slab sups (sup over θ = 0, ½,
1) and the running sup of k\_jm.

\textbf{C3 (windows, item 3).} z\_j = (y∂\_y)ʲv = Σ\_m S(j,m) y\^{}m
w\_m (Stirling), so √E\_j\^{}\{{[}y,1.5y{]}\}{[}err{]} ≤ κ\_ch Σ\_m
S(j,m)√F\_m. For E₃ this is √F₃ + 3√F₂ + √F₁. Here κ\_ch is T9's chart
factor (≤ 1.02).

\textbf{C4 (pointwise level 3, item 4).} For I = {[}y, 1.5y{]}:
sup\_I\textbar w₃\textbar² ≤ ∫\_Iw₃²/\textbar I\textbar{} +
2‖w₃‖\_\{L²(I)\}‖w₄‖\_\{L²(I)\}, with ‖w₃‖² ≤ F₃/(28y⁴) and either ‖w₄‖²
≤ F₄/(28y⁶) or ‖w₄‖² ≤ D₂(τ*)/(28y²).

\textbf{C5 (good time).} Integrating (F) at j = 2 over a window W =
{[}τ\_a, τ\_b{]} with ω ≡ 1 gives θ₂∫\_W D₂ ≤ F₂(τ\_a) + ∫\_W(2g₂√F₂ +
h₂ + inflow). So there is τ* ∈ W with D₂(τ*) ≤ \textbar W\textbar⁻¹∫\_W
D₂. Lemma O does not make this selection: its H4, ∫(h₀⁽⁵⁾)² ≤ M² per
cylinder, is a hypothesis on the data at t₁ = t(τ*), so step 4 must
choose τ* (review A1). The bound is quadratic in Ȳ and is dominated by
F₂(τ\_a)/(θ₂\textbar W\textbar): with the hand-off sub-window W′
(\textbar W′\textbar{} = 1.46·10⁻⁴) it is D₂(τ*) ≤ 1.25·10⁻⁸ at Ȳ =
10⁻¹², and with one whole slab (\textbar W\textbar{} = 0.15) it is
1.26·10⁻¹¹ (\texttt{camp.py\ -\/-window\ 180-180}). On a cylinder
support S (y ≥ y₀ ≥ 22) the localized dissipation obeys ∫\_S w₄² ≤
D\_S/(28y₀²) with D\_S := ∫\_S T₂w₄² ≤ D₂. A localized energy ∫T₂ψ²w₂²
(ψ = 1 on S) does not improve the leading term: its value at τ\_a is
bounded only by F₂(τ\_a), and its inflow (the ramp, or transport through
an inner cut-off) is that of F₂. So the window length, not the
localization, decides. \texttt{p2/t8se/h4loc.py} evaluates H4 per
cylinder (error part L⁷·1.001²·D₂(τ*)/(28y₀²), base part from the T6
hulls of κ̃\^{}\{(≤3)\}). With W′, H4 fails by ×16--279 on the cone
cylinders (r\_c = 200--380) and by ×4.95 on the far cylinder centred at
376. Even the Ȳ-free part, 2.6·10⁻¹⁰ (the review's 2.9·10⁻⁹ also
contains the part linear in Ȳ), fails by ×5.8 at r\_c = 380. With one
slab (slab 179 or 180), the worst cylinder (cone r\_c = 370--380) has
∫(h₀⁽⁵⁾)² ≤ 0.31·M²; the far cylinders have ≤ 0.0052·M²; and the Ȳ-free
part is ≤ 2.7·10⁻¹³, i.e.~≤ 0.6\% of M². Item 4 then uses the same τ*,
and the binding requirement becomes item 3 (Ȳ ≤ 1.79·10⁻¹¹). H1--H4 must
then hold over the whole window slab. By Lemma N3′, level 4 is not
needed for the nonlinear bootstrap either.

\hypertarget{remarks}{%
\subsection{6. Remarks}\label{remarks}}

\textbf{6.1 Nonlinear and modulation terms, levels 0--3 (Lemma N3′):
exact operator, x-dependent majorants, no level 4.}

This replaces the earlier Lemma N3. That lemma expanded ∂⁴(𝒩+𝒬) fully,
which produced the trilinear classes w₃²w₄ and w₃³, and it used T5's box
with placeholder majorants a₀ = 20η\_F, a₁ = 40(η\_F/x +
\textbar κ̃\textbar). On base v2, T7's Hessian sups grow with x (x²n₂ ≤
6.2·10¹⁴ at slab 180), so the placeholders were not justified. N3′ uses
the exact operator, has no trilinear classes, covers every level j ≤ 3
and the modulation part 𝒬, and is evaluated on the certified far boxes.

\emph{Operator.} Let α := 1 − κ̃e and J² := α² + e'². At the base
modulation and with the clock fixed (T7 \texttt{nl7.py}), the equation
is e\_τ = F(e) with

F(e) = (J/α){[}−κ\_ss + c(−q̃e' + (p̃+e)α)/J + b̃(−t̃₁e' + ñα)/J{]} − Ṽ +
e'(ϖ − eθ̃\_τ)/α.

F = A₄e'\,'\,'\,' + A₃e'\,'\,' + A₂ with A₄ = −J⁻⁴ and A₃, A₂ functions
of z := (e, e', e'\,'). The nonlinear remainder is

𝒩 := F(e) − F(0) − DF(0)e = Ãe'\,'\,'\,' + B̃e'\,'\,' + C̃, Ã := 1 − J⁻⁴,
B̃ := A₃(z) − A₃(0), C̃ := A₂(z) − A₂(0) − ∇A₂(0)·z.

\emph{Cancellation: no coefficient grows with x.} Insert the base
identities

ϖ = −κ̃κ̃' + Q + c(q̃ − x) + b̃t̃₁ + A, θ̃\_τ = −κ̃'\,'\,' − κ̃²κ̃' + κ̃Q − cxκ̃ +
ρ\_θ with ρ\_θ = r' + κ̃A, Ṽ = −κ̃'\,' + cp̃ + b̃ñ + r.

The quadratic part of the tangential term is ee'(κ̃ϖ − θ̃\_τ)/α = ee'(cκ̃q̃
+ b̃κ̃t̃₁ + κ̃'\,'\,' − r')/α. Its first two parts cancel the quadratic
parts −cq̃κ̃ee'/α and −b̃t̃₁κ̃ee'/α of the normal velocity. The remaining c-
and b̃-terms are linear in e. Hence

𝒩 = 𝒩\_K + (κ̃'\,'\,' − r')ee'/α, where 𝒩\_K is the nonlinear part of
−(J/α)κ\_ss.

So 𝒩 depends on none of c, b̃, p̃, q̃, ñ, t̃₁, x, Q and A. Its coefficients
are functions of κ̃, \ldots, κ̃'\,'\,' and r' only, and none of them grows
with x. \texttt{n3sym.py} checks A₄ = −J⁻⁴, F(0) = −r and the
independence of c, \ldots, Q symbolically; A drops out of the generated
forms. T7's growth of n₂ comes from the interval hull of the cancelling
sum, with q̃, x and Q taken as independent intervals; it is not a
property of 𝒩.

\emph{Modulation part.} For the terms 𝒬, with the true clock and
modulation (T5 §2.1: c → (1−ϑ)c − x₁, b̃ → (1−ϑ)b̃ − x₂, κ\_ss →
(1−ϑ)κ\_ss), the exact split is

𝒩 + 𝒬 = (1−ϑ)𝒩 + ϑe'\,'\,'\,' + ϑκ̃²e'\,' − ϑc₀e + ϑP + g·e' − (ϑc +
x₁)e,

P := (κ̃'\,'\,' − r')ee'/α, g := U/α + 5ϑκ̃κ̃`, U := (x₁ + ϑc)q̃ + (x₂ +
ϑb̃)t̃₁, c₀ = −3κ̃'² − 4κ̃κ̃''.

\texttt{n3qsym.py} checks this split symbolically. The first five terms
have the structure of 𝒩, with à → (1−ϑ)à + ϑ and C̃ → (1−ϑ)C̃ + ϑκ̃²e'\,' −
ϑc₀e + ϑP, and no coefficient grows with x; they go through the weak
form below. The last two terms carry the velocity U ≈ (x₁ + ϑc)x, which
does grow with x; they go through the strong form. The modulation bounds
come from T5's (M): \textbar x\textbar{} ≤ m̄(1 − m̄ē(A+W₁))⁻¹(ℓ̄A + c\_φℛ₂
+ n\_φη₁(A+W₁+B) + strip). The linear part is T7's C\_xΛ\^{}\{1/2\}, and
the residual part x\_r := m̄c\_φℛ₂ is added explicitly per slab (T7's m̄;
c\_φ is the 2-norm of (c\_φ,j)\_\{j=1..3\}, the largest over the T7
certificates, since the largest component understates it ×1.30 (T7);
T6's ℛ₂; x\_r ≤ 2.6·10⁻¹⁶ on slabs 2--180), in camp and here (review
B6). The nonlinear and E\_ij parts are O(η₁Ȳ) relative. T7 Lemma 4c(M)
bounds them on the box by x\_NL ≤ 1.95·10⁻¹⁵. This fits in the slack
C\_x(Ȳ − Y\_cert) ≥ 2.44·10⁻¹² (ratio ≤ 8.0·10⁻⁴; t7/constants.json
meta.modulation\_nl) that the certified (Y) level leaves for every Ȳ ≥
10⁻¹². That level is 1.33·10⁻¹³ (p2/ycomp.py, SLAB0 = 2), with Y(τ₀) ≤
6.21·10⁻¹³ at the switch. So \textbar x\textbar{} ≤ C\_xȲ + x\_r and
\textbar ϑ\textbar{} = \textbar x₃\textbar/n ≤ (C\_xȲ + x\_r)/n, with
the per-slab lower bound of n = ‖χṼ‖\_ρ (cconst's nV).

\emph{Weak form, level 3.} Recall w\_m = ∂\^{}\{m+1\}e. At level 3, 𝒩
contributes S := 2∫T₃w₃∂⁴𝒩 to F₃'. Two integrations by parts give S =
2∫(T₃w₃)'\,'∂²𝒩, with

∂²𝒩 = Ãw₅ + a₄w₄ + a₃w₃ + a♭,

a₄ := 2Ã' + B̃, a₃ := Ã'\,' + 2B̃' + (∂\_\{e'`\}A₃)e''\,' + ∂\_\{e'`\}C̃,
a♭ := {[}B̃'\,'{]}♭e''\,' + {[}C̃'\,'{]}♭.

Here ' is the total x-derivative along the solution, and {[}·{]}♭ drops
the chain-rule term that contains e'\,'\,'`. Since à depends on (e, e')
and A₃, A₂ depend on z, the coefficients a₄, a₃ and those inside a♭
involve only e, \ldots, e'\,'`, the boxed jets. So ∂²𝒩 is linear in the
free jets w₃, w₄, w₅, and \textbf{there are no trilinear classes}. Write
T := T₃ and ϕ := T'(6Ã' + 4B̃) − 2T'\,'Ã. One more integration by parts
of the w₄w₅- and w₄²-terms gives

S = 2∫TÃw₅² + ∫Ψ₄₅w₄w₅ + ∫Ψ₃₅w₃w₅ + ∫Ψ₃₄w₃w₄ + ∫Ψ₃₃w₃² + 2∫(Tw₅ + 2T'w₄
+ T'\,'w₃)a♭,

Ψ₄₅ = 2Ta₄, Ψ₃₅ = 2Ta₃ + 4T'`Ã − T'(6Ã' + 4B̃), Ψ₃₄ = T'(4a₃ − 6Ã'\,' −
4B̃`) − 2T'`B̃ + 2T'`'Ã, Ψ₃₃ = 2T''a₃.

There are no boundary terms. At y\_a, T, \ldots, T'\,'\,' vanish (W2).
At ℓ̃ the closed-curve argument of (SE) applies, because 𝒩 is even when e
is. S involves w\_\{≤5\}, κ̃\^{}\{(≤5)\} and r\^{}\{(≤3)\} only, so the
identity extends by density as in §1.5.

\emph{Weak form, levels j ≤ 2.} For j = 1, 2 the contribution
2∫T\_jw\_j∂\^{}\{j+1\}𝒩 equals 2∫(T\_jw\_j)'\,'∂\^{}\{j−1\}𝒩, with
∂\^{}\{j−1\}𝒩 = Ãw\_\{j+2\} + α₁w\_\{j+1\} + α₀w\_j + α♭ and

j = 2: α₁ = Ã' + B̃, α₀ = B̃' + ∂\_\{e'`\}C̃, α♭ = ∂\_XC̃ + (∂\_eC̃)e' +
(∂\_\{e'\}C̃)e'`; j = 1: α₁ = B̃, α₀ = ∫₀¹∂\_\{e'`\}C̃(e, e', se'`)ds, α♭ =
C̃(e, e', 0).

With ϕ := 4T'α₁ − 2T'`Ã − 2T'Ã', the same steps give S\_j =
2∫TÃw\_\{j+2\}² + Σ∫Ψ\_\{ab\}w\_aw\_b + 2∫(Tw\_\{j+2\} + 2T'w\_\{j+1\} +
T'\,'w\_j)α♭, where

Ψ\_\{j+1,j+2\} = 2Tα₁, Ψ\_\{j,j+2\} = 2Tα₀ + 4T'`Ã − T'(4α₁ − 2Ã'),
Ψ\_\{j,j+1\} = T'(4α₀ − 4α₁' + 2Ã'`) − 2T'`α₁ + 4T'`Ã' + 2T'`'Ã,
Ψ\_\{jj\} = 2T''α₀

(j = 3 is the case above). At level 0, 2∫T₀w₀∂𝒩 = −2∫(T₀w₀)'𝒩, and one
more integration by parts of the Ãw₃-terms gives

S₀ = 2∫TÃw₂² + ∫2T(Ã' − B̃)w₁w₂ + ∫{[}4T'`Ã + 2T'(2Ã' − B̃){]}w₀w₂ +
∫{[}2T'`'Ã + 4T'`Ã' + 2T'Ã'\,'{]}w₀w₁ − 2∫(Tw₁ + T'w₀)C̃,

with \textbar C̃\textbar{} ≤ γ\_e\textbar e\textbar{} +
γ₀\textbar w₀\textbar{} + γ₁\textbar w₁\textbar, γᵢ := ½ΣⱼZⱼ
sup\textbar∂ᵢ∂ⱼA₂\textbar. \texttt{n3wcheck.py} checks every one of
these identities (the difference of the integrands is an exact
derivative; Euler-operator test).

\emph{Strong form for g·e' and (ϑc + x₁)e.} At level j,
2∫T\_jw\_j∂\^{}\{j+1\}(gw₀) = ∫{[}(2j+1)g'T − gT'{]}w\_j² +
2Σ\_\{k=2\}\^{}\{j+1\} C(j+1,k)∫Tw\_j g\^{}\{(k)\}w\_\{j+1−k\}. The flux
Tgw\_j² vanishes at ℓ̃ because g(ℓ̃) = 0 (q̃(ℓ̃) = t̃₁(ℓ̃) = κ̃`(ℓ̃) = 0 on base
v2), and T = 0 at y\_a. The x-growth of U is absorbed: gτ₁ is bounded
because τ₁ = O(1/Y\_ℓ). In g\^{}\{(j+1)\} the top jet e\^{}\{(j+1)\} =
w\_j occurs once, with coefficient ∂\_eg = Uκ̃/α²; that part costs 2
sup\textbar∂\_eg·e'\textbar{} in σ\_j. Hence

Δσ\_j ∋ sup\_C{[}(2j+1)\textbar g'\textbar{} +
\textbar g\textbar\textbar τ₁\textbar{} +
2\textbar∂\_eg\textbar\textbar e'\textbar{]} + 2(x̄ +
ϑ̄\textbar c\textbar) (at j = 0 the ∂\_eg-term is part of g'), Δk\_jj ∋
sup\_R{[}((2j+1)\textbar g'\textbar{} +
2\textbar∂\_eg\textbar\textbar e'\textbar)T +
\textbar g\textbar\textbar T'\textbar{]}/W, Δc\_\{j,j+1−k\} = 2C(j+1,k)
sup\_\{C∪R\}\textbar g\^{}\{(k)\}♭\textbar Y\_ℓ\^{}\{k−1\} (k = 2,
\ldots, j+1; ♭: the top jet removed),

where x̄ := C\_xȲ + x\_r and ϑ̄ := x̄/n.~The Δc\_\{jm\} enter g\_j as
½Δc\_\{jm\}√F\_m. The majorants of g\^{}\{(k)\} use the Frenet jets of
q̃, t̃₁ (q̃' = 1 + κ̃p̃, p̃' = −κ̃q̃, t̃₁' = κ̃ñ, ñ' = −κ̃t̃₁) and (1/α)\^{}\{(m)\},
m ≤ 4, from \texttt{n3qforms.py}.

\emph{Boxes.} The bootstrap hypothesis 𝔹′ is the following, for the slab
and x ≥ x\_a:

\begin{itemize}
\tightlist
\item
  y ≥ 28: \textbar e\textbar{} ≤ 1.7·E(x) with E := η₁ +
  √F₀·y\^{}\{3/2\}/√84 and √F₀ ≤ A₀Ȳ + B₀ (camp's split, per slab), and
  \textbar e'\textbar{} ≤ ζ₁, x\textbar e'`\textbar{} ≤ ζ₂,
  x²\textbar e''\,'\textbar{} ≤ ζ₃;
\item
  y \textless{} 28 (x ≤ 34.25): \textbar e\textbar,
  \textbar e'\textbar{} ≤ η₁ = 10⁻⁵ (T5's tip box, T5 §2.5; T7 Lemma 4b
  certifies it: c\_e1Ȳ = 2.8·10⁻⁸ at Ȳ = 1.7·10⁻¹², and
  \texttt{n3boxes.py} checks this), and x\textbar e'`\textbar,
  x²\textbar e''`\textbar{} ≤ η\_F = 10⁻² (the e'' part is the
  assembly's Lemma 4a).
\end{itemize}

The defaults are ζ₁ = 1.5·10⁻³, ζ₂ = 3·10⁻⁴, ζ₃ = η\_F, with Ȳ =
1.7·10⁻¹² in E and in the modulation bound. \texttt{n3boxes.py} checks
all four for the final camp run, at Ȳ = 1.7·10⁻¹² (every smaller Ȳ is
then covered, since the bounds are monotone in Ȳ): \textbar e'\textbar{}
≤ 5.47·10⁻⁴, x\textbar e'`\textbar{} ≤ 1.20·10⁻⁴,
x²\textbar e''\,'\textbar{} ≤ 9.33·10⁻³, and √F₀ (with N3′) ≤ 1.00·√F₀
(used). At Ȳ = 10⁻¹² the margins are ×2.76, ×3.97, ×1.69 and ×1.70
(review B2: the first version assumed 7·10⁻⁴, 1.65·10⁻⁴ and factor 1.25,
so its true margins were ×1.29, ×2.19, ×1.76 and ×1.25). The
\textbar e\textbar{} box is a fixed point: n3far is run with the camp
output without N3′, then once more with the camp output of that n3 file.
(E bounds \textbar e\textbar{} because \textbar e\textbar{} ≤
\textbar e(28)\textbar{} + H₂₈√F₀, with \textbar e(28)\textbar{} ≤ η₁
and H₂₈(y)² := ∫₂₈\^{}y T₀⁻¹ ≤ (y³ − 28³)/84.)

\emph{Majorants per x-range.} Cut each slab's far field into the ramp R,
{[}x\_b, 30.25{]}, {[}30.25, 34.25{]} and the dyadic ranges {[}2ᵏ·34.25,
2ᵏ⁺¹·34.25{]} up to ℓ̃. On each range, the base data are arb hulls of κ̃,
\ldots, κ̃\^{}\{(5)\}, r', r'`, r''\,' and p̃, q̃, ñ, t̃₁ over the range and
the slab. The box enters through (Z₀, Z₁, Z₂, Z₃) := (1.25E(x\_hi), ζ₁,
ζ₂/x\_lo, ζ₃/x\_lo²) for y ≥ 28, and (η₁, η₁, η\_F/x\_lo, η\_F/x\_lo²)
below. The box is convex and contains 0, so mean-value forms with sups
over the box give (i, j = 0, 1, 2 index z):

\begin{itemize}
\tightlist
\item
  Ã⁺, \textbar Ã\textbar, \textbar Ã'\textbar{} and
  \textbar Ã'\,'\textbar{} directly;
\item
  \textbar B̃\textbar{} ≤ Σᵢ Zᵢ sup\textbar∂ᵢA₃\textbar,
  \textbar B̃`\textbar{} ≤ Σᵢ Zᵢ sup\textbar∂\_X∂ᵢA₃\textbar{} + Σᵢ Zᵢ₊₁
  sup\textbar∂ᵢA₃\textbar, \textbar∂\_\{e'`\}C̃\textbar{} ≤ Σⱼ Zⱼ
  sup\textbar∂ⱼ∂\_\{e'`\}A₂\textbar; then \textbar a₄\textbar{} ≤
  2\textbar Ã'\textbar{} + \textbar B̃\textbar{} and
  \textbar a₃\textbar{} ≤ \textbar Ã'`\textbar{} +
  2\textbar B̃'\textbar{} + Z₃ sup\textbar∂\_\{e'`\}A₃\textbar{} +
  \textbar∂\_\{e''\}C̃\textbar;
\item
  \textbar a♭\textbar{} ≤ β♭\_e\textbar e\textbar{} + Σ\_\{m≤2\}
  β♭\_m\textbar w\_m\textbar. In {[}C̃'\,'{]}♭ = ∂\_X²C̃ + 2Σᵢ(∂\_X∂ᵢC̃)zᵢ'
  + Σᵢⱼ(∂ᵢ∂ⱼA₂)zᵢ'zⱼ' + Σ\_\{i≤1\}(∂ᵢC̃)zᵢ'\,' and in
  {[}B̃'\,'{]}♭e'\,'\,', every monomial is one jet times boxed factors.
  Keep the jet and box the rest. The α♭ of the levels j ≤ 2 and the γ's
  of level 0 are treated the same way.
\end{itemize}

For 𝒬, all majorants of 𝒩 are multiplied by 1 + ϑ̄, Ã⁺ and
\textbar Ã\textbar{} gain ϑ̄(1 + \textbar Ã\textbar) and ϑ̄,
\textbar∂\_\{e'`\}C̃\textbar{} gains ϑ̄κ̃², the Hessian sups of A₂ gain ϑ̄
times those of P, and the linear terms ϑκ̃²e'' and −ϑc₀e add their own
derivatives to the α₀ and β♭ of each level. ∂\_X differentiates through
the base only (κ̃\^{}\{(k)\} → κ̃\^{}\{(k+1)\}, r\^{}\{(k)\} →
r\^{}\{(k+1)\}). \texttt{n3forms.py} (generated by \texttt{n3sym.py})
and \texttt{n3qforms.py} (generated by \texttt{n3qsym.py}) give every
derivative needed as an arb function.

\emph{Constants.} On the core (y ≥ y\_b), with τ\_m := T\^{}\{(m)\}/T,
set

m₃₅ := 2\textbar a₃\textbar{} + 4\textbar τ₂\textbar\textbar Ã\textbar{}
+ \textbar τ₁\textbar(6\textbar Ã'\textbar{} + 4\textbar B̃\textbar), m₃₄
:= \textbar τ₁\textbar(2\textbar Ã'`\textbar{} + 4\textbar B̃'\textbar{}
+ 4Z₃ sup\textbar∂\_\{e'`\}A₃\textbar{} + 4\textbar∂\_\{e''\}C̃\textbar)
+ 2\textbar τ₂\textbar\textbar B̃\textbar{} +
2\textbar τ₃\textbar\textbar Ã\textbar,

m₄₅ := 2\textbar a₄\textbar, m₃₃ :=
2\textbar τ₂\textbar\textbar a₃\textbar,

so that \textbar Ψ\_\{ab\}\textbar{} ≤ m\_\{ab\}T (sups over the core
pieces; m₄₅ also over R₃). On the ramp set M₃₅ := sup\_R Ψ₃₅²/(TW), and
M₃₄ likewise, using the same bounds with T\^{}\{(m)\} in place of Tτ\_m
(T\^{}\{(m)2\}/T is bounded by (W2)). Also set M₃₃ := sup\_R
\textbar Ψ₃₃\textbar/W, M\_f := sup\_R T'`²/(TW), M\_T := sup\_R
\textbar T'`\textbar/W and τ̄₂ := sup\_C \textbar τ₂\textbar. Young's
inequality is used with the weights δ\_D = 10, δ\_R = 10⁻², δ\_A = 0.05,
δ\_B = 3·10⁻³, δ\_C = 1, δ\_f = 10⁻³ and δ₃₄ := (4M₃₄/M\_T)\^{}\{1/2\},
together with N₄² ≤ D₃ + ¼F₃ + ½∫(T'')⁺w₃² (S3). With c\_N := ½m₃₄ +
m₄₅/(2δ\_D) + ½δ₃₄ + δ\_B this gives

Δθ₃ = 2 sup Ã⁺ + ½(m₃₅ + m₄₅)δ\_D + δ\_R + δ\_A + c\_N,

Δσ₃ = m₃₅/(2δ\_D) + ½m₃₄ + m₃₃ + δ\_C τ̄₂² + c\_N(¼ + ½τ̄₂),

Δk₃₃ = M₃₅/(4δ\_R) + M₃₄/(2δ₃₄) + M₃₃ + δ\_f M\_f + ½c\_N M\_T.

The a♭ terms are paired as follows, with ‖a‖²\_T := ∫Ta² and
‖a‖²\_\{(1)\} := ∫(T'²/T)a² (the script doubles the 1/δ\_A and 4/δ\_B
terms, which is spare):

2∫Tw₅a♭ ≤ δ\_A D₃ + ‖a♭‖²\_T/δ\_A, 4∫T'w₄a♭ ≤ δ\_B N₄² +
4‖a♭‖²\_\{(1)\}/δ\_B,

2∫T'\,'w₃a♭ ≤ δ\_C τ̄₂²F₃ + δ\_f M\_f V₃² + ‖a♭‖²\_\{T,C\}/δ\_C +
‖a♭‖²\_\{T,R\}/δ\_f.

Three facts are used: (Σ of four terms)² ≤ 4Σ(·)²; T₃/T\_m =
Y\_ℓ\^{}\{6−2m\} on the core and y\^{}\{6−2m\} on the ramp (one ramp for
all levels); and on y ≥ 28, e² ≤ 2η₁² + 2H₂₈²F₀. Together they add to h₃
the terms Σ\_\{m≤2\}(Δd₃ₘF\_m + Δk₃ₘV\_m²) + hf₃, where

Δd₃ₘ = 4{[}f\_C sup\_C β♭\_m²Y\_ℓ\^{}\{6−2m\} + f₁ sup\_C
τ₁²β♭\_m²Y\_ℓ\^{}\{6−2m\} + f\_R sup\_R β♭\_m²y\^{}\{6−2m\}{]}, plus
8{[}f\_C∫\_C Tβ♭\_e²H₂₈² + f₁∫\_C Tτ₁²β♭\_e²H₂₈²{]} for m = 0,

Δk₃ₘ = 4f₁ sup\_R β♭\_m²T'²/(TW),

hf₃ = 8η₁²{[}f\_C∫\_C Tβ♭\_e² + f₁∫\_C Tτ₁²β♭\_e²{]} + 4η₁²{[}f\_R∫\_R
Tβ♭\_e² + f₁∫\_R(T'²/T)β♭\_e²{]},

f\_C = 2/δ\_A + 1/δ\_C, f\_R = 2/δ\_A + 1/δ\_f, f₁ = 8/δ\_B.

Levels j = 1, 2 use the same formulas with (j, j+1, j+2) in place of (3,
4, 5), T\_j/T\_m = Y\_ℓ\^{}\{2(j−m)\}, j + 1 terms in α♭ instead of 4,
and the weights δ\_D = 1, δ\_R = δ\_A = 10⁻³, δ\_B = 10⁻⁴; θ₂ enters
item 4 through C5, so the D-side weights are kept small there. Level 0
pairs the three Ψ's as above (with N₁, √F₀, √D₀ and the ramp datum V₀),
and the C̃-term as 2∫T\textbar w₁\textbar\textbar C̃\textbar{} ≤ 2γ₁N₁² +
2γ₀√F₀N₁ + 2N₁(E\_abs + E\_F√F₀), and so on for T'w₀C̃; here E\_abs² :=
2η₁²∫\_C Tγ\_e² + η₁²∫\_R Tγ\_e² and E\_F² := 2∫\_\{y≥28\}Tγ\_e²H₂₈².
Every product of two norms is split by Young (weights 10⁻⁴ on N₁² for
the absolute and ramp T'-terms), and N₁² ≤ D₀ + (¼ + ½τ̄₂)F₀ + ½M\_T V₀²
(S3).

\textbf{Lemma N3′.} Assume 𝔹′, (SE), \textbar x\textbar{} ≤ x̄ and
\textbar ϑ\textbar{} ≤ x̄/n.~Then (F) at each j ≤ 3 holds for the
equation with 𝒩 + 𝒬 included, provided θ\_j, λ\_j, k\_jj, k\_jm and
c\_jm are replaced by θ\_j − Δθ\_j, λ\_j − Δσ\_j, k\_jj + Δk\_jj, k\_jm
+ Δk\_jm and c\_jm + Δc\_jm, and h\_j gains
Σ\_\{m\textless j\}Δd\_jmF\_m + hf\_j. All of these are per-slab arb
bounds.

\emph{Evaluation (base v2, slabs 2--180; certified).} The inputs are the
SE-corrected T6 data that \texttt{p2/se\textbackslash{}\_corr.py} writes
(p2/t6/jets and residuals.json), with its outward rounding. Run
\texttt{cconst.py} (TOPLOSS = 0.0196) and \texttt{merge.py} to get
constants\_fin.json, and \texttt{camp.py} without N3′ to get
amp\_fin.json. Then run
\texttt{n3far.py\ -\/-slabs\ 2-180\ -\/-amp\ constants/amp\textbackslash{}\_fin.json\ -\/-out\ constants/n3\textbackslash{}\_it0.json}
(about 30 s, one process) and
\texttt{camp.py\ ...\ -\/-n3\ constants/n3\textbackslash{}\_it0.json}.
Then run n3far once more with \texttt{-\/-amp} set to that camp output
(→ n3.json), camp again, and \texttt{n3boxes.py}. Every output JSON has
an \texttt{args} block with the argv, the env and the input paths. The
cconst slab files also record sha256 digests of their jets, of
residuals.json and of the T7 inputs. A cached slab is reused only if
these match, and merge.py checks that they agree. camp enters each
Δd\_jmF\_m through √F\_m ≤ A\_mȲ + B\_m, with an adaptive split of the
cross term.

\begin{itemize}
\tightlist
\item
  sup Ã⁺ = 0.0161 at slab 180 on the 1.7-box (sup\textbar κ̃e\textbar{} ≤
  4.1·10⁻³; the certified value is 2.4·10⁻³). The top loss is 2 sup Ã⁺ ≈
  8 sup\textbar κ̃e\textbar. The assembly's TOPLOSS = 2 sup\textbar A₄ −
  A₄(0)\textbar{} ≤ 0.0196, on the unscaled box, is the same quantity.
\item
  Level 3: Δθ₃ ≤ 0.096, which takes θ₃ from 0.220 to ≥ 0.124; Δσ₃ ≤
  2.9·10⁻³; Δk₃₃ ≤ 1.7·10⁹ (+0.15\% of k₃₃); Δd₃ₘ ≤ (1.6·10⁵, 74,
  0.044); Δk₃ₘ ≤ (9.5·10³, 8.0·10⁵, 4.0·10⁷); Δc₃ₘ ≤ (0.044, 2.8·10⁻³,
  9.3·10⁻⁵); hf₃ ≤ 9.1·10⁻¹².
\item
  Levels 0--2: Δθ\_j ≤ 0.035 (0.032 of it is the top term); Δσ\_j ≤
  2.0·10⁻³, 1.3·10⁻³, 1.3·10⁻³; Δk\_jj ≤ 7.3, 8.3·10³, 4.6·10⁶ (≤ 0.05\%
  of k\_jj); Δd₁₀ ≤ 1.4, Δd₂ₘ ≤ (2.5·10³, 0.18); Δc\_jm ≤ 7.1·10⁻⁴;
  hf\_j ≤ 3·10⁻¹³.
\item
  𝒬: x̄ ≤ 1.33·10⁻¹¹ and ϑ̄ ≤ 5.8·10⁻⁸ (at Ȳ = 1.7·10⁻¹²; slab 180, n ≥
  2.3·10⁻⁴). Its strong-form share of Δσ\_j is ≤ 1.7·10⁻⁶; it is the
  only source of the Δc\_jm.
\item
  Requirement (pipeline of STEP4\_RESULT, current KCH, review B4 and B6
  applied). On W′ (SUBW): Ȳ ≤ 4.99·10⁻¹² (item 4), against 5.11·10⁻¹²
  without N3′. On a one-slab window (\texttt{-\/-window\ 180-180}): Ȳ ≤
  1.79·10⁻¹¹ (item 3, E3), against 1.80·10⁻¹¹ without N3′.
\item
  𝔹′ closes, as checked above (thinnest margin ×1.69,
  x²\textbar e'\,'`\textbar). N3′ is certified for Ȳ ≤ 1.7·10⁻¹² (the
  run's Ȳ); the x²\textbar e''\,'\textbar{} box itself closes up to Ȳ ≈
  1.83·10⁻¹².
\item
  \emph{Review B1: the residual of the corrected base, computed
  directly.} \texttt{p2/t6c/slabc.py} runs T6's engine on the
  (SE)-corrected base itself, so se\_corr's majorants are no longer
  used. a₁ = −δ is exact. a₂ is the root of Im Ŵ(ℓ) = 0; on every slab,
  by Banach, it is unique in a ball of radius ≤ 7.6·10⁻²⁶ (contraction ≤
  2·10⁻¹⁴). a₁′ and a₂′ come from T6's θ\_τ = ρ\_θ + Φ − Φ(0) and from
  its velocity identity. The geometry uses e\^{}\{iθ\} from θ itself, at
  time degree 20. The two window elements are cut at u = 0.6 and 0.8. On
  slab 11 a knot crosses u = 0.6, so that slab is also cut in time. The
  corrected far end encloses 0: \textbar Im Ŵ(ℓ)\textbar{} ≤ 2.2·10⁻²⁰
  and \textbar θ̂(ℓ) + 3π/2\textbar{} ≤ 1.9·10⁻²⁵. With the bump set to
  zero, the run reproduces T6. Maxima over slabs 2--180, direct against
  se\_corr: R1 4.6·10⁻¹⁵ vs 2.6·10⁻¹⁴; 𝓡₀ 3.5·10⁻⁹ vs 8.7·10⁻⁹; 𝓡₁
  4.4·10⁻¹¹ vs 4.5·10⁻¹⁰; 𝓡₂ 2.5·10⁻¹⁰ vs 6.9·10⁻⁹; 𝓡₃ 4.9·10⁻⁹ vs
  1.7·10⁻⁷. se\_corr left out that S₁(ℓ) moves; this shifts ∂⁻¹r over
  the whole far field. The direct level-0 hats are up to ×3 per slab,
  but their maxima are below se\_corr's. Chain on these data
  (\texttt{constants/*\textbackslash{}\_b1*}): B\_j = 7.6·10⁻⁷,
  1.8·10⁻⁹, 5.1·10⁻⁸, 7.8·10⁻⁶. The far-box margins at Ȳ = 10⁻¹² are
  ×6.43, ×4.69, ×1.99 and ×1.70. The requirement is Ȳ ≤ 5.15·10⁻¹² on W′
  and ≤ 1.79·10⁻¹¹ on one slab. H4 on one slab is ≤ 0.23·M². ϱ\_c is
  still 0.04711. The (Y) level is 3.2·10⁻¹⁴.
\end{itemize}

\emph{Remarks.}

\begin{itemize}
\tightlist
\item
  The class terms (A), (B) of the old N3 in \texttt{cconst.py} (κ\_N,
  \texttt{NLFAC}) are superseded; keeping them is conservative (θ₃
  −5·10⁻⁵ and k₃₃ +2.7·10⁵, approximately).
\item
  With \texttt{-\/-n3}, cconst's TOPLOSS counts the top term a second
  time; it may be set to 0.
\item
  \texttt{n3far.py\ -\/-noQ\ -\/-levels\ 3} reproduces the level-3-only
  evaluation of the first version of N3′ (with that version's boxes).
\end{itemize}

\emph{Conclusion.} N3′ uses F\_j, D\_j, N\_\{j+1\} (through S3), the
lower F\_m and the ramp data V₀--V₃. \textbf{Level 4 is not used.}
Together with C5, level 4 is dropped altogether.

\textbf{6.2 Variants.}

\begin{itemize}
\tightlist
\item
  Nested ramps (R\_j ⊂ C\_\{j−2\}) replace W-data at level j by
  D\_\{j−2\}. This is good for level 4 (A₄ falls 40×) and bad for level
  3.
\item
  Narrow or steep ramps lose through (S1): the P²/T term at the steep
  foot of S dominates the transport flux.
\item
  With δ\_P = 1.5 (most of D\_j used) the best level-3 ramp is {[}14,
  22{]}.
\end{itemize}

\textbf{6.3 Localization.} A transport-following outer factor
Ω(Y\_ℓ/Y\_h), with Y\_h' = c\_hY\_h and c\_h \textless{} c, removes the
bend from the hand-off energies. However:

\begin{itemize}
\tightlist
\item
  its 4th-order commutators at small Y\_h cost more than the bend
  contributes;
\item
  in the global T5 design the bend contributes ≤ 3\% of √F₃, through
  modulation (γ\_j on the bend), and \textless{} 10⁻⁵ through c\_ji.
\end{itemize}

It is therefore not recommended.

\textbf{6.4 Not covered.} These remain outside the lemma:

\begin{itemize}
\tightlist
\item
  the T7 constants: Lemma J inputs, C\_x and the modulation constants
  are taken from t7/constants.json, some still pending or provisional
  there; μ\_e, ϱ\_c and K\_e are placeholders;
\item
  T9's chart factors and the base's own window energies (T1);
\item
  hypothesis (SE) at the far end (§3.1, Remark), and T4's F\_j(τ₀);
\item
  the bootstrap boxes of 6.1 at Ȳ above 1.7·10⁻¹² (they close up to Ȳ ≈
  1.83·10⁻¹²), and the tip box η₁ of 6.1 (review A3).
\end{itemize}

\end{document}
